\section{Square presentations and their translation periods}\label{square-presentations-and-their-translation-periods}

Let
\[
 a(x,y)=(-x,y),\quad b(x,y)=(y,x),\quad c(x,y)=(x,1-y).
\]
\textbf{Proposition UA.Affine (the complete square presentation).} These generate \(W=\mathbb Z^2\rtimes D_8\), of type \(\{4,4\}\), with
\(ac=ca\). Its presentation is
\[
 W=\langle a,b,c\mid a^2=b^2=c^2=(ab)^4=(bc)^4=(ac)^2=1\rangle. \tag{S1}
\]
\emph{Proof.} Put \(d=bab\), \(y_0=cd\), and
\(x_0=by_0b\). The relations give \([a,d]=1\), \([a,y_0]=1\),
\(dy_0d=y_0^{-1}\), \(ax_0a=x_0^{-1}\), and conjugation by \(b\) swaps
\(x_0,y_0\). Substitution of \(bc=x_0ab\) makes \((bc)^4=1\) equivalent
to \([x_0,y_0]=1\). Thus every word collects to \(x_0^u y_0^v M\),
\(M\in\langle a,b\rangle=D_8\). In the affine model \(x_0,y_0\) are the
unit translations, so these forms are distinct. This proves (S1). \(\square\)

The normal form also fixes our affine multiplication convention: the rightmost transformation acts first. The symbol \(D_8\) denotes the group of signed permutation matrices in two coordinates, of order eight. In all dihedral notation below the subscript denotes the order of the group.

\textbf{Proposition UA.Square (finite square seeds).} For \(n\ge5\), define the invariant lattice
\[
 \Lambda_n=\begin{cases}
  2^e\mathbb Z^2,&n=2e+3,\\
  2^e\langle(1,1),(1,-1)\rangle,&n=2e+4,
 \end{cases}\qquad e\ge1,
\]
and put \(H_n=W/\Lambda_n\). Its order is \(2^n\); the full linear
\(D_8\) coordinate is retained, even if its action on the vertex quotient
is not faithful. Its pair orders are \(4,4,2\). It has all Boolean deletions and IP. Its third lower-central subgroup is \(2\mathbb Z^2/\Lambda_n\), of order \(2^{n-5}\).

\emph{Proof.} The lattice indices are respectively \(2^{2e}\) and \(2^{2e+1}\), so the normal form gives the stated orders. The three linear pair products have orders four, four and two; hence the marked pair orders cannot fall in either lattice quotient.

Write \(T=abcb\), \(U=cbab\), translations by \((-1,0)\), \((0,1)\).
The extra relators for \(H_n\) can be taken as \(T^{2^e},U^{2^e}\) in
the square case, and \(T^{2^e}U^{2^e},T^{2^e}U^{-2^e}\) in the diagonal
case. After deleting any marker, both \(T,U\) have order at most two:
deleting \(a\) leaves conjugates of \(c\), deleting \(c\) leaves
conjugates of \(a\), and deleting \(b\) leaves \(ac,ca\). Therefore
every lattice relator survives every deletion. The even Coxeter relators
also survive. Thus \(H_n\) has Boolean deletions and IP by UA.D.

Moreover,
\[
 \gamma_3(W)=2\mathbb Z^2,
 \qquad \gamma_3(H_n)=2\mathbb Z^2/\Lambda_n,
 \qquad |\gamma_3(H_n)|=2^{n-5}.                         \tag{S2}
\]
Indeed reduction modulo \(2\mathbb Z^2\) makes the adjacent squares
central; conversely \([[a,b],c]\) and \([[b,c],a]\) give the two doubled
unit translations, up to signs. Thus \(H_5\) is class two of order 32.
The squared translations are not always independent cyclic subgroups:
in \(H_6\) they coincide and have order two. Formula (S2) uses a lattice
index and does not require independence of those cyclic subgroups. \(\square\)
